\section{Evaluation}
\label{sec:eval}

We implement \sys{} in CuTe DSL (CUTLASS~4.8, CUDA~13.0) and evaluate it on NVIDIA H100 (80~GB HBM3) and B200 (180~GB HBM3e) GPUs with $H_Q=32$, $H_{KV}=8$, $D=128$, page size 64, and token-wise K and V quantization groups of $g=128$ values with FP32 scales and offsets.
Token-wise K and V exercise both acc-fold on QK and operand-fold on PV; \Cref{sec:appendix-axes} evaluates all four combinations of token-wise and channel-wise K and V, where \sys{} achieves $1.75$--$4.00\times$ speedups over FlashInfer BF16 on H100 and $1.50$--$2.74\times$ on B200.
We evaluate batch sizes 1 and 8 at context lengths of 8K, 32K, 128K, and 512K, and batch size 64 at 8K, 16K, 32K, 64K, and 128K.
Each measurement takes the median of five timed iterations (excluding an L2 flush before each call) after five warmups, and includes all decode kernels and split reduction.
Plotted speedups normalize each kernel's latency to FlashInfer BF16, whereas speedups quoted in the text compare \sys{} directly against each named baseline within the same benchmark series (using five-run medians for B200 INT4 vendor baselines and three-run medians for BitDecoding and SAW).

\subsection{Kernel performance}
\label{sec:eval-performance}

\begin{figure}[t]
  \centering
  \captionsetup[subfloat]{farskip=2pt}
  \includegraphics[width=\columnwidth]{assets/speedup-subbyte-legend.pdf}
  \par
  \subfloat[H100, batch 1.\label{fig:speedup-subbyte-h100-b1}]{%
    \includegraphics{assets/speedup-subbyte-h100-b1.pdf}}
  \hfill
  \subfloat[H100, batch 8.\label{fig:speedup-subbyte-h100-b8}]{%
    \includegraphics{assets/speedup-subbyte-h100-b8.pdf}}
  \hfill
  \subfloat[H100, batch 64.\label{fig:speedup-subbyte-h100-b64}]{%
    \includegraphics{assets/speedup-subbyte-h100-b64.pdf}}
  \par\vspace{3pt}
  \subfloat[B200, batch 1.\label{fig:speedup-subbyte-b200-b1}]{%
    \includegraphics{assets/speedup-subbyte-b200-b1.pdf}}
  \hfill
  \subfloat[B200, batch 8.\label{fig:speedup-subbyte-b200-b8}]{%
    \includegraphics{assets/speedup-subbyte-b200-b8.pdf}}
  \hfill
  \subfloat[B200, batch 64.\label{fig:speedup-subbyte-b200-b64}]{%
    \includegraphics{assets/speedup-subbyte-b200-b64.pdf}}
  \caption{Kernel speedups of \sys{} and 2- and 4-bit baselines over FlashInfer BF16.}
  \label{fig:speedup-subbyte}
\end{figure}

\begin{figure}[t]
  \centering
  \captionsetup[subfloat]{farskip=2pt}
  \includegraphics[width=\columnwidth]{assets/speedup-wide-legend.pdf}
  \par
  \subfloat[H100, batch 1.\label{fig:speedup-wide-h100-b1}]{%
    \includegraphics{assets/speedup-wide-h100-b1.pdf}}
  \hfill
  \subfloat[H100, batch 8.\label{fig:speedup-wide-h100-b8}]{%
    \includegraphics{assets/speedup-wide-h100-b8.pdf}}
  \hfill
  \subfloat[H100, batch 64.\label{fig:speedup-wide-h100-b64}]{%
    \includegraphics{assets/speedup-wide-h100-b64.pdf}}
  \par\vspace{3pt}
  \subfloat[B200, batch 1.\label{fig:speedup-wide-b200-b1}]{%
    \includegraphics{assets/speedup-wide-b200-b1.pdf}}
  \hfill
  \subfloat[B200, batch 8.\label{fig:speedup-wide-b200-b8}]{%
    \includegraphics{assets/speedup-wide-b200-b8.pdf}}
  \hfill
  \subfloat[B200, batch 64.\label{fig:speedup-wide-b200-b64}]{%
    \includegraphics{assets/speedup-wide-b200-b64.pdf}}
  \caption{Kernel speedups of \sys{} and 8- and 16-bit baselines over FlashInfer BF16.
  The dashed line marks parity.}
  \label{fig:speedup-wide}
\end{figure}


\Cref{fig:speedup-subbyte,fig:speedup-wide} compare \sys{} against sub-byte and wider-format baselines.
On H100, \sys{} INT2 achieves a $3.36\times$ geometric-mean speedup over FlashInfer BF16 and $2.03\times$ over TensorRT-LLM FP8, with maxima of $4.15\times$ and $2.44\times$.
On B200, the geometric-mean speedups are $2.67\times$ over BF16, $1.48\times$ over TensorRT-LLM FP8, and $1.16\times$ over TensorRT-LLM NVFP4, with maxima of $3.25\times$, $1.64\times$, and $1.26\times$.
These vendor kernels are invoked through FlashInfer, and \sys{} INT2 outperforms them across all tested batch sizes and context lengths.
Against OSCAR INT2 and BitDecoding INT2, \sys{} INT2 achieves geometric-mean speedups of $2.31\times$ and $1.77\times$ on H100, and $3.44\times$ and $2.78\times$ on B200.

On H100, \sys{} INT4 achieves geometric-mean speedups of $2.83\times$ over FlashInfer BF16, $1.71\times$ over TensorRT-LLM FP8, $1.32\times$ over BitDecoding INT4, and $4.19\times$ over SAW INT4.
On B200, \sys{} INT4 achieves geometric-mean speedups of $2.30\times$ over FlashInfer BF16, $1.28\times$ over TensorRT-LLM FP8, $1.94\times$ over BitDecoding INT4, and $6.93\times$ over SAW INT4.
It matches native NVFP4 overall ($1.00\times$, ranging from $0.88\times$ to $1.20\times$; faster at 32K--128K for batch 64), while returning BF16 output instead of NVFP4's FP8 output.
\Cref{sec:appendix-head-scaling} further evaluates how \sys{} INT2 speedup scales across $H_{KV}\in[1,32]$.

\begin{figure}[t]
  \centering
  \captionsetup[subfloat]{farskip=2pt,justification=raggedright,singlelinecheck=false}
  \includegraphics[width=\columnwidth]{assets/fold-ab-grid-legend.pdf}
  \par
  \begin{minipage}[t]{0.345\columnwidth}
  \centering
  \subfloat[H100: \mbox{reconstruction} versus folding.\label{fig:fold-ab-hopper}]{%
    \includegraphics[width=\linewidth]{assets/fold-ab-grid.pdf}}
  \end{minipage}\hfill
  \begin{minipage}[t]{0.307\columnwidth}
  \centering
  \subfloat[B200: E4M3 \mbox{reconstruction} versus folding.\label{fig:fold-ab-b200-folding}\label{fig:fold-ab-grid-b200}]{%
    \includegraphics[width=\linewidth]{assets/fold-ab-grid-b200-folding.pdf}}
  \end{minipage}\hfill
  \begin{minipage}[t]{0.307\columnwidth}
  \centering
  \subfloat[B200: folded E4M3 versus E2M1.\label{fig:fold-ab-b200-feed}]{%
    \includegraphics[width=\linewidth]{assets/fold-ab-grid-b200-feed.pdf}}
  \end{minipage}
  \caption{\sys{} INT2 kernel latency ratios between ablation variants at equal split-KV counts (first variant's latency divided by the second's).
  Panels (b) and (c) decode an INT2 cache into E4M3 or E2M1 Tensor Core operands on B200; error bars show paired 95\% confidence intervals across ten runs.}
  \label{fig:fold-ab-grid}
\end{figure}


\subsection{Kernel accuracy}
\label{sec:eval-accuracy}

\begin{figure}[t]
  \centering
  \includegraphics[width=\columnwidth]{assets/accuracy-arms-legend.pdf}
  \par
  \subfloat[H100.\label{fig:accuracy-arms-h100}]{%
    \includegraphics[width=0.49\columnwidth]{assets/accuracy-arms-h100.pdf}}
  \hfill
  \subfloat[B200.\label{fig:accuracy-arms-b200}]{%
    \includegraphics[width=0.49\columnwidth]{assets/accuracy-arms-b200.pdf}}
  \caption{$L_2$ output error of BF16 decode (FlashInfer), \sys{} INT4, \sys{} INT2, and OSCAR INT2 on the synthetic input (logarithmic scale).}
  \label{fig:accuracy-arms}
\end{figure}


\Cref{fig:accuracy-arms} compares the $L_2$ output error of BF16 decode (FlashInfer), \sys{} INT4, \sys{} INT2, and OSCAR INT2.
To stress-test quantization, our synthetic input amplifies eight K channels, ramps K's per-token scale along the sequence, and rescales K so that pre-softmax attention logits have unit standard deviation.
For INT2, \sys{} and OSCAR decode the same quantized cache written by OSCAR's clipped quantizer without BF16 windows.
Both INT2 kernels yield $L_2$ errors of $0.52$--$0.61$, compared with $0.16$--$0.22$ for \sys{} INT4 and $(2.13\text{--}2.25)\times10^{-3}$ for BF16 decode.
Under unclipped min--max INT2 quantization, the $L_2$ error rises to $0.86$--$1.03$.
Compared directly against the dequantized INT2 cache, \sys{}'s kernel arithmetic error is only $(1.68\text{--}1.76)\times10^{-3}$ on both H100 and B200, matching BF16 decode and confirming that the $0.52$--$0.61$ error stems from 2-bit quantization rather than affine folding.

\subsection{Kernel ablation}
\label{sec:eval-ablation}

On INT2 caches with equal split-KV counts within each pair, \Cref{fig:fold-ab-grid} ablates affine folding against per-value reconstruction on H100 and B200, and FP4 (E2M1) against FP8 (E4M3) KV Tensor Core operands on B200.
On H100, \sys{}'s speedup over external baselines combines its warp-specialized kernel design with affine folding, which adds a $1.07$--$1.15\times$ speedup over per-value reconstruction (\Cref{fig:fold-ab-hopper}) as expected when H100 reconstruction remains within the hardware issue peak (\Cref{fig:ops-per-byte}).
On B200, the reconstruction baseline follows the standard approach of dequantizing INT2 values into FP8 (E4M3) Tensor Core operands before each matrix product; because this per-value reconstruction exceeds B200's issue peak, affine folding achieves a $7.36$--$8.86\times$ speedup (\Cref{fig:fold-ab-b200-folding}).
Switching the folded B200 INT2 kernel from FP8 (E4M3) to FP4 (E2M1) KV operands and their matching packing layout provides an additional $1.17$--$1.21\times$ speedup (\Cref{fig:fold-ab-b200-feed}).

\subsection{Serving throughput}
\label{sec:eval-serving-throughput}

\begin{figure}[t]
  \centering
  \captionsetup[subfloat]{farskip=2pt}
  \includegraphics[width=\columnwidth]{assets/serving-pareto-h100-legend.pdf}
  \par
  \subfloat[H100, 99\% prefix-cache hit.\label{fig:serving-pareto-h100-99}]{%
    \includegraphics{assets/serving-pareto-h100-99.pdf}}
  \hfill
  \subfloat[H100, 95\% prefix-cache hit.\label{fig:serving-pareto-h100-95}]{%
    \includegraphics{assets/serving-pareto-h100-95.pdf}}
  \par\vspace{3pt}
  \subfloat[B200, 99\% prefix-cache hit.\label{fig:serving-pareto-b200-99}]{%
    \includegraphics{assets/serving-pareto-b200-99.pdf}}
  \hfill
  \subfloat[B200, 95\% prefix-cache hit.\label{fig:serving-pareto-b200-95}]{%
    \includegraphics{assets/serving-pareto-b200-95.pdf}}
  \caption{Per-user and aggregate Qwen3-8B serving throughput on a single GPU with 100K input and 1K output tokens.}
  \label{fig:serving-pareto}
\end{figure}


We serve Qwen3-8B with 100K input tokens, 1K output tokens, and 99\% or 95\% prefix-cache hits (one cached prefix per request), sharing a BF16 prefill backend across configurations.
We evaluate two INT2 configurations to show both realistic accuracy-preserving integration and pure 2-bit kernel speed: \emph{OSCAR INT2 + Ours} replaces OSCAR's INT2 decode kernel with \sys{} while keeping OSCAR's rotations and BF16 windows (128 initial and 512 recent tokens, $<0.7\%$ of the 100K context), whereas \emph{Ours INT2} uses \sys{}'s token-wise INT2 K and V cache without recent windows or rotations.
We vary concurrency in powers of two up to 16 on H100 and 32 (plus 48) on B200 to measure per-user (median request output tokens/s) and aggregate (total output tokens/s) throughput (\Cref{fig:serving-pareto}).

At a 99\% prefix-cache hit rate, \emph{OSCAR INT2 + Ours} reaches a peak aggregate throughput of 487 tokens/s on H100 and 868 tokens/s on B200, compared with 269 and 333 tokens/s for OSCAR ($81\%$ and $161\%$ higher).
\emph{Ours INT2} reaches 506 tokens/s on H100 and 897 tokens/s on B200, improving over OSCAR by $88\%$ and $170\%$; this also exceeds FlashInfer BF16's peaks of 115 tokens/s on H100 (at concurrency 2) and 355 tokens/s on B200 (at concurrency 8) by $341\%$ and $153\%$.
At 95\% prefix hits, \emph{OSCAR INT2 + Ours} and \emph{Ours INT2} reach 337 and 346 tokens/s versus OSCAR's 214 tokens/s on H100, and 738 and 755 tokens/s versus 297 tokens/s on B200.
Across cache formats, peak throughput reflects both decoder speed and the number of concurrent requests that fit in memory.

\subsection{Serving accuracy}
\label{sec:eval-serving-accuracy}

\begin{table}[t]
  \centering
  \caption{Qwen3-8B serving accuracy on H100 (mean $\pm$ sample standard deviation across five runs).}
  \label{tab:serving-accuracy}
  \footnotesize
  \setlength{\tabcolsep}{4pt}
  \begin{tabular}{@{}lrr@{}}
    \toprule
    Configuration & GPQA-Diamond (\%) & GSM8K (\%) \\
    \midrule
    FlashInfer BF16 & $56.36 \pm 1.84$ & $95.71 \pm 0.11$ \\
    OSCAR INT2 & $54.65 \pm 2.18$ & $95.74 \pm 0.38$ \\
    OSCAR INT2 + Ours & $55.45 \pm 1.49$ & $95.92 \pm 0.15$ \\
    \bottomrule
  \end{tabular}
\end{table}


\Cref{tab:serving-accuracy} compares FlashInfer BF16, OSCAR INT2, and \emph{OSCAR INT2 + Ours} on GPQA-Diamond~\cite{rein2023gpqa} and GSM8K on H100, reporting the mean and sample standard deviation across five runs per configuration.
Every question-paired 95\% confidence interval for the accuracy difference between configurations spans zero, so five runs show no measurable accuracy gain or loss (though five runs do not establish formal equivalence).
Median prompt lengths are 243 tokens for GPQA-Diamond and 91 tokens for GSM8K; because \emph{OSCAR INT2 + Ours} decodes the identical quantized cache as OSCAR with $<1.8\times10^{-3}$ arithmetic error (\Cref{sec:eval-accuracy}), it preserves OSCAR's cache accuracy while accelerating decode.
